% Required preamble commands/packages:
% \usepackage{amsmath,amssymb,amsthm}
% \newtheorem{corollary}{Corollary}
% \newcommand{\E}{\mathbb{E}}
% \newcommand{\R}{\mathbb{R}}
% \newcommand{\diag}{\operatorname{diag}}
% \newcommand{\rhoop}{\rho}

\section{Proof of the Damping Corollary}
\label{app:damping_corollary}

\begin{corollary}[Slow-subspace damping]
Let the true slow dynamics have eigenvalue $\lambda_s^{\mathrm{true}}$, and
let the learned model's slow imagination dynamics have eigenvalue
$\lambda_s^{\mathrm{model}}$ along the same one-dimensional slow direction.
Assume $\lambda_s^{\mathrm{model}}>\lambda_s^{\mathrm{true}}>0$ and let $U$
be an orthonormal basis aligned with the true slow drift subspace. Then damping
the model update along $U$ with rate $\eta$ reduces slow-subspace one-step MSE
if and only if
\begin{equation}
    0 < \eta
    <
    2\left(1 -
    \frac{\lambda_s^{\mathrm{true}}}{\lambda_s^{\mathrm{model}}}
    \right).
\end{equation}
\end{corollary}

\begin{proof}
Consider first the scalar slow coordinate. Let $s_t$ denote the coordinate of
the latent state in the true slow direction. The true and model transitions are
\begin{equation}
    s_{t+1}^{\mathrm{true}}
    =
    \lambda_s^{\mathrm{true}} s_t,
    \qquad
    s_{t+1}^{\mathrm{model}}
    =
    \lambda_s^{\mathrm{model}} s_t.
\end{equation}
The undamped one-step slow-subspace error is
\begin{equation}
    e_0
    =
    \left(
        \lambda_s^{\mathrm{model}}
        -
        \lambda_s^{\mathrm{true}}
    \right)s_t.
\end{equation}

Damping along the aligned slow direction scales the model's slow update by
$(1-\eta)$, so the damped transition is
\begin{equation}
    s_{t+1}^{\eta}
    =
    (1-\eta)\lambda_s^{\mathrm{model}}s_t.
\end{equation}
Equivalently, in vector form the effective damped transition is
\begin{equation}
    A_\eta = A - \eta U U^\top A,
    \label{eq:damped_transition}
\end{equation}
and on the aligned slow subspace its spectral radius is
\begin{equation}
    \rhoop(A_\eta|_{\mathcal{S}_{\mathrm{true}}})
    =
    |1-\eta|\,|\lambda_s^{\mathrm{model}}|.
\end{equation}
For the usual damping range $0 \leq \eta \leq 1$, this is
$(1-\eta)|\lambda_s^{\mathrm{model}}|$, so any positive damping decreases the
model's slow spectral radius before over-correction is considered.
TODO: If the final paper uses damping of the update
$\Delta x_t=(A-I)x_t$ rather than damping of $A x_t$, replace
\eqref{eq:damped_transition} with
$A_\eta=A-\eta U U^\top(A-I)$ and update the scalar bound accordingly.

The damped one-step slow-subspace error is
\begin{equation}
    e_\eta
    =
    \left(
        (1-\eta)\lambda_s^{\mathrm{model}}
        -
        \lambda_s^{\mathrm{true}}
    \right)s_t.
\end{equation}
Since $s_t^2 \geq 0$, damped MSE is lower than undamped MSE if and only if
\begin{equation}
    \left|
        (1-\eta)\lambda_s^{\mathrm{model}}
        -
        \lambda_s^{\mathrm{true}}
    \right|
    <
    \left|
        \lambda_s^{\mathrm{model}}
        -
        \lambda_s^{\mathrm{true}}
    \right|.
    \label{eq:damping_error_inequality}
\end{equation}
Let
\begin{equation}
    a=\lambda_s^{\mathrm{model}},
    \qquad
    b=\lambda_s^{\mathrm{true}},
\end{equation}
with $a>b>0$. Then \eqref{eq:damping_error_inequality} becomes
\begin{equation}
    |(1-\eta)a-b| < |a-b|.
\end{equation}
Because $a-b>0$, this is equivalent to
\begin{equation}
    -(a-b) < (1-\eta)a-b < a-b.
\end{equation}
The right inequality gives
\begin{equation}
    (1-\eta)a-b < a-b
    \quad
    \Longleftrightarrow
    -\eta a < 0
    \quad
    \Longleftrightarrow
    \eta > 0.
\end{equation}
The left inequality gives
\begin{align}
    -(a-b) &< (1-\eta)a-b \\
    -a+b &< a-\eta a-b \\
    0 &< 2a - \eta a - 2b \\
    \eta &< 2\left(1-\frac{b}{a}\right).
\end{align}
Substituting back $a=\lambda_s^{\mathrm{model}}$ and
$b=\lambda_s^{\mathrm{true}}$ yields
\begin{equation}
    0 < \eta
    <
    2\left(1 -
    \frac{\lambda_s^{\mathrm{true}}}{\lambda_s^{\mathrm{model}}}
    \right).
\end{equation}

The upper bound is the over-correction threshold. For small positive $\eta$,
the damped model eigenvalue moves toward the true slow eigenvalue. If $\eta$ is
too large, the damped transition crosses past the true transition and the
slow-subspace error grows again, potentially with a sign flip in the slow
coordinate.

For an $r$-dimensional aligned slow subspace, the same argument applies
component-wise after diagonalizing the slow block. If the slow block is
$\Lambda_s^{\mathrm{model}}=\diag(\lambda_1^{\mathrm{model}},\ldots,\lambda_r^{\mathrm{model}})$
and the true block is
$\Lambda_s^{\mathrm{true}}=\diag(\lambda_1^{\mathrm{true}},\ldots,\lambda_r^{\mathrm{true}})$,
then a shared scalar damping rate must satisfy the bound for every slow mode:
\begin{equation}
    0 < \eta <
    \min_i 2\left(
        1-\frac{\lambda_i^{\mathrm{true}}}{\lambda_i^{\mathrm{model}}}
    \right).
\end{equation}
TODO: Add the non-diagonal slow-block version using operator norms or a
Lyapunov/MSE comparison in the slow coordinates.
\end{proof}
